% ============================================================
%  R. Nielsen, "Om det oprindelige Forhold mellem Religion og Videnskab"
%  Kjøbenhavn: J. H. Schultz, 1881
%  Indbydelsesskrift til Kjøbenhavns Universitets Fest, 8. april 1881
%  TRANSCRIPTION (Danish) — from Royal Danish Library digital scan
% ============================================================
\documentclass[12pt,a4paper]{article}
\usepackage[utf8]{inputenc}
\usepackage[T1]{fontenc}
\usepackage[danish]{babel}
\usepackage{microtype}
\usepackage{setspace}
\usepackage[margin=1.2in]{geometry}
\usepackage{fancyhdr}
\usepackage{hyperref}

\hypersetup{
  pdftitle={Om det oprindelige Forhold mellem Religion og Videnskab},
  pdfauthor={R. Nielsen},
  hidelinks
}

\pagestyle{fancy}
\fancyhf{}
\rhead{Nielsen, \emph{Om det oprindelige Forhold} (1881)}
\cfoot{\thepage}

\setstretch{1.3}

\title{\textbf{Om det oprindelige Forhold\\[4pt]
       mellem Religion og Videnskab}}
\author{R.\ Nielsen}
\date{Kjøbenhavn: J.\ H.\ Schultz, 1881}

\begin{document}
\maketitle
\thispagestyle{fancy}

\bigskip

\noindent Den bittre Strid mellem Religion og Videnskab vil ikke lade
sig hæve, med mindre det skeer paa den ufravigelige Betingelse,
at begges oprindelige Selvstændighed gjensidigt anerkjendes. Religionen
for sig maa da kunne bestaae uden al Videnskab, og Videnskaben for sig
uden al Religion; men hvorledes er det muligt? At Religionen ikke er
af videnskabelig Oprindelse, »den aabenbarede Religion« saa lidt som
Mytherne, maa vel indrømmes; men Historien lærer, at Videnskaben neppe
slaaer Øinene op, før den begynder at rokke ved den givne Religion og
indlede Striden. Hos Grækerne var det især Philosopherne, som førte
Angrebet. Xenophanes paaviste Modsigelser i den polytheistiske
Folketro; imod de mange Guder satte han Eenheden, imod deres timelige
Oprindelse Evigheden, mod det Ubestandige i deres Væsen
Uforanderligheden, imod deres Lighed med Mennesker Ophøietheden
o.\ s.\ v.\footnote{Zeller, \emph{Die Philosophie der Griechen}, Ister
Theil, 2te Auflage, Tubingen 1856, s.\ 381.} Hos de store Digtere
sporer man en trinvis Forandring i Gudernes Væsen; hos Euripides begynde
de allerede at hjemfalde til Reflexionen. Den naturlige Forstaaeise
rykker den Dannede nærmere. Kunsten holdt paa Formen, men netop den
plastiske Formfuldendelse tærede paa Gudernes Realitet. Zeus i Olympia
stod vistnok som Mægler mellem Folketroen og den philosophisk-kritiske
Opfattelse; men Indtrykket var æsthetisk, og det overvældende æsthetiske
Indtryk kvalte det Religiøse. Man kan uden Overdrivelse paastaae, at
den videnskabelige Tænkning, som undergravede de græske Guders
Autoritet, havde en mægtig Forbundsfælle i den plastiske Kunst.

Hvad Forholdet mellem det Gamle og Nye Testamentes aabenbarede Religion
angaaer, stiller Sagen sig anderledes. Begavet med rige Evner, dyb
Følelse, storartet Phantasi, skarp Forstand og stærke Lidenskaber, var
Israels Folk dog ikke som Grækerne særlig anlagt enten paa plastisk
Kunst eller dybgaaende videnskabelig Forskning; det var i
Verdenshistorien Jehovas Folk, i særegen Betydning Religionens Folk,
Aabenbaringens udvalgte Folk. En kunstnerisk Efterligning af det
Guddommelige var forbudt, Poesien stod udelukkende i Religionens
Tjeneste, og den kritisk prøvende Videnskab havde ingen Rod i
Folkenaturen. At Moseloven med sin høiere Ethik ikke tilsteder nogen
grundig videnskabelig Undersøgelse, er indlysende. Hvor meget de ti
Bud end synes at tiltale den menneskelige Fornuft, saa er det dog strax
iøinefaldende, at det ikke er Fornuftens Grunde, men Jehovas Ord,
hvorpaa Eftertrykket falder. Man lægge vel Mærke til, at den samme
Gud, der aabenbarer sin Villie i den almindelige Lov, ogsaa udtaler sig
directe i de særlige Love.

Angaaende Tabernaklet siger Jehova (2den Moseb.\ XXV): »Og dette er
Offergaven, som I skulle tage af dem: Guld og Sølv og Kobber; Mørkeblaat
og Purpurrødt og Carmosinrødt og fint Linned og Gedehaar« o.\ s.\ v.
»Og Tabernaklet (XXVI) skal du gjøre af ti Forhæng \ldots\ Længden af
hvert Forhæng skal være otte og tyve Alen og Breden fire Alen paa hvert
Forhæng, eet Maal for alle Forhæng« \ldots\ Ingen skal sige, at disse
specielle Bud overskride Fornuftens Fatteevne; det Overnaturlige, for
os Ufattelige ligger tvertimod deri, at den Almægtige, Himlens og
Jordens Skaber, har talt saa nøiagtigt om Tepper og Sløifer og heri
slet Intet overladt til Moses' sunde Sands. Det er dette, som her maa
forekomme Fornuften aldeles ubegribeligt. Paa modsat Maade ydmyges
Fornuften ikke blot i det Gamle, men ogsaa i det Nye Testamente, naar
det Overnaturlige slaaer igjennem saa pludselig, saa vældig, at Naturen
revner, Lovenes Baand brister og Miraklet skeer. Eller staaer maaskee
Begivenheden med Gadarenernes Sviin ikke lige saa høit over Fornuften,
som Miraklet (2den Kong.\ XIII) med den døde Mand, der faldt ned i
Propheten Elisas Grav: »der kom Liv i ham, og han reiste sig op paa
sine Been«. Der viser sig i slige Træk en saadan Splid mellem den
guddommelige Aabenbaring og den menneskelige Fornuft, at ethvert
Forlig mellem Religion og Videnskab maa synes umuligt. Skal
Aabenbaringen gjælde, saa maa Videnskaben bort; skal Videnskaben
gjælde, saa maa Aabenbaringen opløses: et fortvivlet Dilemma.

Vi ville her ikke opholde os ved enkelte Træk af de Stridigheder, som
nu i næsten to tusinde Aar ere førte mellem Aabenbaringstroen og
Fornuften, Stridigheder, hvorunder snart Religionen har løftet det
kirkelige Banner høit og været i Færd med at underkue Videnskaben,
snart igjen den seirende Videnskab paa den krænkede Fornufts Vegne
været i Begreb med at opløse Religionen; vor Opgave er at gaae tilbage
til Stridens oprindelige Grund, for at undersøge, om der i Grunden
selv ikke skulde skjule sig en Misforstaaelse, en uklar Sammenblanding
og Forvexling af religiøs og videnskabelig Erkjendelse. Christi Ord
(Joh.\ VIII, 32): »I skulle erkjende Sandheden og Sandheden skal
frigjøre Eder«, hæve sig høit over alle Partistridigheder, og maae da
lige saa fuldt komme Videnskaben som Religionen til gode.

Men, indvender man, den Sandhed, der skal frigjøre os, maa nødvendig
stemme overeens med sig selv i Et og Alt. Hvis Noget, seet fra en
Side, vel er sandt, fra en anden Side derimod usandt: da ligger dette
ikke i Sandheden, men i dens Fremstilling og Opfattelse. Er end al vor
Erkjendelse Stykværk, saa kan det til Sandheden selv svarende
Erkjendelsesprincip dog umulig lade sig udstykke. To indbyrdes
stridende Erkjendelsesprinciper maae jo forudsætte to principielt
modsatte, mod hinanden stridende Sandheder, men det er aldeles
urimeligt. Dersom den religiøse Erkjendelse staaer i aabenbar Strid
med den videnskabelige, da kan dette ikke betyde, at her virkelig er
to Sandheder og for Erkjendelsen to absolut uensartede Principer; men
et af de to foregivne Principer maa nødvendig være falsk, enten det
religiøse eller det videnskabelige. Har det religiøse Princip Gyldighed,
saa maa det omfatte hele Sandheden, ogsaa den videnskabelige; Videnskaben
kan da ikke gjøre Fordring paa at have noget særeget Princip. Er det
omvendt det videnskabelige Princip, der har Gyldighed, da er Religionen
i sig selv principløs, og maa i Et og Alt være Videnskabens Dom
underlagt. --- Vi skulle i det Følgende undersøge, hvorledes det hermed
forholder sig.

At Sandheden i evig Forstand maa være een og udeelbar, ville vistnok
Alle indrømme; men Sandhedens Væsen er i sig selv dobbeltsidig,
hverken blot noget Objectivt ei heller noget blot Subjectivt, men en
Erkjendelsens Eenhed af begge. Uden erkjendende Subject gives der ingen
Erkjendelse, og uden et Object, som bliver erkjendt, gives der heller
ingen Erkjendelse. Sandheden selv kan altsaa hverken være det
Subjective for sig ei heller det Objective for sig; den gjennemtrængende
Eenhed af Subjectivt og Objectivt maa gaae igjennem Altilværelsen,
omfattende den hele Virkelighed med alle dens Muligheder i det Mindste
som i det Største. Dette maa Mennesket med sin ufuldkomne Erkjendelse
indrømme; men hermed indrømmer det tillige, at det i sin Begrændsning
umulig kan gjennemskue den absolute Sandheds Indhold i dettes utallige
Former og Concretioner. Skal Mennesket virkelig erkjende Sandheden, da
maa Sandheden dele sig; og den deler sig som Lyset. Idet Sandhedens
Lys straaler ind i Menneskets Bevidsthed, afsætter det to
Erkjendelses-Midtpunkter, to Principer, som hvert for sig antage sit
relative Præg af det Absolute, nemlig Troens og Videns; det første
angaaer væsenlig Menneskets eget Selv, dets Indre, og er personligt;
det andet angaaer Tingene og deres væsenlige Sammenhæng, og er for saa
vidt upersonligt.

Skiller man Tro og Viden fra de tilsvarende Principer, vilde de
bestandig flyde over i hinanden. At Troen umulig kan undvære al Viden,
er lige saa indlysende som at Viden ikke kan undvære al Tro. Den
Troende maa jo nødvendig vide, hvad det er, hvorpaa han troer og hvad
det at troe vil sige; lige saa maa den Vidende troe paa sine Sandser og
sin Forstand. Skeptikeren, der hverken kan troe paa det Ene eller det
Andet, ei en Gang paa sin egen Skepticisme, efterdi Sandheden er, at
der ingen Sandhed er, er i lige Grad afspærret fra Videnskaben og fra
Religionen.

Der er en dybt gaaende Forskjel mellem religiøs og videnskabelig Tro,
mellem religiøs og videnskabelig Viden. Den religiøse Tro har et
psychisk-ethisk Præg og er væsenlig en Samvittighedssag, et reent
personligt Anliggende; den videnskabelige Tro angaaer ikke Mennesket
særligt, men udtrykker Tillid til Videnskabens Grundbetingelser og
Realitet. Den religiøse Tro sætter Kundskab om Gud og det Overnaturlige
langt høiere end Kjendskab til Virkelighedens Love og de naturlige
Ting: al tilsvarende Viden faaer herved alene sin Betydning. Det
Overnaturlige kan ikke sandses, og maa dog give sig tilkjende paa
umiddelbar Maade; det antager da en sandselig-oversandselig Form i
Phantasien, og gjør derigjennem Indtryk paa Følelserne: det aabenbarer
sig. Athene aabenbarede sig for Achileus (Il.\ I, 194--98), Jehova
aabenbarer sig for Abraham, for Moses, for Propheterne; selv Guds
Aabenbaring i Christus gjør herfra ingen Undtagelse. I Johannes første
Brev I, 1 hedder det vel: »Det, som var fra Begyndelsen, det, som vi
have hørt, det, som vi have seet med vore Øine, det, som vi have
beskuet, og vore Hænder have følt paa, nemlig om Livets Ord«; det
kunde vel synes, som blev det Overnaturlige her ligefrem bekræftet af
Sandserne og bevidnet paa sandselig Maade; men sammenligne vi hermed
(Matth.\ XVI, 15--17), hvor Jesus siger til Peder, der erkjender ham
for Guds Søn, »salig er du Simon, Jonas's Søn, thi det har Kjød og
Blod ikke aabenbaret dig, men min Fader, som er i Himlen«, saa komme
vi til fuld Forvisning om, at en sand Viden om det Overnaturlige maa
være af overnaturlig Oprindelse. Den Viden, der har den religiøse Tro
til Grundlag, kan altsaa ikke opløse Troen; naar Troen svinder, blive
de overnaturlige Objecter øieblikkeligt forvandlede til Myther og
Indbildninger, Jehova og Christus lige saa vel som Zeus og Athene. I
denne Henseende har Anselmus med sit \emph{credo, ut intelligam}
unegtelig Ret imod Abailard, der vilde vende Forholdet om.

Anderledes i Videnskaben. Vel kan man ogsaa sige, at den videnskabelige
Viden begynder med Tro; men deels er Visheden om det for Bevidstheden
Nærværende ikke Tro, men en umiddelbar Viden; deels er Alt hvad der
foreløbig antages i Tro, fordi det antages i uprøvet Tillid, bestemt
til at underkastes en Prøvelse, hvorved Troen, ifølge det Cartesianske:
\emph{de omnibus dubitandum est}, opløser sig i Tvivl, og naar saa
Tvivlen er hævet ved Begrundelsen, forvandler sig til Viden. Vi begynde
med at orientere os i Verden; Sandser, Forstand og Fornuft virke
sammen paa umiddelbar Maade; snart er det Sandserne, snart Forstanden,
der har eensidig Overvægt, medens Fornuften dog altid kræver, at
Vexelbestemmelsen af det, som sandses, og det, som tænkes, holdes i
Hævd, thi herpaa beroer den videnskabelige Viden. Idet Undersøgelsen
begynder angaaende Forskjellen mellem hvad der i Tilværelsen er
Foranderligt og hvad der som et Uforanderligt maa danne Grundlaget for
alle Forandringer, tager Videnskaben sin Begyndelse. Opgaven er, at
undersøge Forholdet mellem Tingene og Lovene. Tingene sandses, Lovene
tænkes; men da Tingene lige saa lidt kunne existere uden Love, som
Lovene kunne tænkes uden Ting, saa maa Sandsning og Tænkning oprindelig
forudsætte hinanden i Fornuftens Eenhed.

At den religiøse og den videnskabelige Erkjendelse maae ifølge den
angivne Karakteristik være væsenlig forskjellige, er da vel indlysende.
Den religiøse Tro, der har sin Viden i Phantasi og Selvfølelse, kaster
et Blik paa Tingenes naturlige Verden, forvisser sig om dens
Forkrænkelighed og griber det Overnaturlige; tynget af Loven vender
den sig mod Miraklerne og seer i dem Vidnesbyrd om, at den frie,
almægtige Gud er Lovenes Herre. De religiøse Love ere Frihedslove,
Opdragelseslove, der kunne forandres efter Tid og Omstændigheder,
saasom Lovene om rene og urene Dyr; de videnskabelige Love,
mathematiske, mechaniske, physisk-chemiske o.\ s.\ v., høre derimod
saa væsenlig med i Tingenes Natur, at de umulig kunne forandres, uden
Naturen selv maa forandres. Alle Videnskabslove ere Nødvendighedslove;
kunde de brydes eller ophæves, saa maatte den hele Verden falde sammen;
Videnskaben maa desaarsag i Fornuftens Navn resolut angribe det
Mirakuløse.

Der er vel dem, der mene, at dette er en for haard Tale. Der er i den
virkelige Tilværelse ganske vist Knudepunkter, hvori det Overnaturlige
og det Naturlige, Miraklet og Loven synes at brydes med hinanden; men
det er kun tilsyneladende: det nye Reale, der indsættes ved Miraklet,
maa jo ved at indgaae i Sammenhæng med det Naturlige, rette sin
Udvikling efter Loven. Vistnok, svare vi, kan man i Almeensætningers
svævende Form forlige det indbyrdes Modstridende, men hvad hjælper det,
naar Modsigelsen viser sig i Kjendsgjerningen selv. Der berettes (2den
Moseb.\ III, 2--4) om Mose Kaldelse: »Jehovas Engel lod sig tilsyne
for ham i en Ildslue midt ud af en Tornebusk, og han saae, og see
Tornebusken brændte i Ilden, og Tornebusken blev ikke fortæret. Da
sagde Moses: jeg vil vige til Siden og see det store Syn, hvorfor
Tornebusken ikke opbrændes. Men Jehova saae, at han veg tilside for at
see, og Gud kaldte paa ham midt ud af Tornebusken og sagde: Moses!
Moses! og han sagde: see, her er jeg«. Hvorledes det vil være Troen
muligt at opfatte denne Beskrivelse, kunne vi endda nok forstaae, thi
Troen seer Alt i Phantasiens Lys, og Phantasien veed, hvorledes det
Overnaturlige forbinder sig med det Fornuftige, Lovenes Gyldighed med
Lovenes Afbrydelse; men i Viden og Videnskab lader det sig ikke
opfatte. Den menneskelige Viden er i Alt, baade Sandseligt og
Tænkeligt, saa strengt bundet til det Naturlige, det blot Fornuftige,
at den hverken kan høre Gud tale som et Menneske ei heller indse,
hvorledes en Tornebusk kan vedblive at brænde uden at fortæres af
Ilden. Heraf er det da klart, at et Troens Object maa være af en ganske
anden Natur end et Videns Object. Stille vi to Naturforskere, een paa
hver Side af Moses, de vilde da kunne see Tornebusken, saa sandt det er
en udvortes virkelig Gjenstand, men de vilde umulig enten kunne see
Engelen eller høre Jehovas Ord. Tornebusken er et Videns Object, men
Engelen og Jehova ere Troens Objecter.

Hvad kan nu heraf udledes? At Videnserkjendelsen, Erkjendelsen af det
Naturlige, maa være fælles for alle naturlig Begavede, der gjøre
fornuftig Brug af Sands og Forstand, medens Troeserkjendelsen,
Erkjendelsen af det Overnaturlige, saa sandt en saadan Erkjendelse
virkelig gives, kun er mulig for »de Udvalgte«, for dem, som med særlig
Sands for det Overnaturlige kunne see, hvad Andre ikke kunne see, og
høre, hvad Andre ikke kunne høre.

% -------------------------------------------------------
% Section II onwards — transcription forthcoming
% -------------------------------------------------------

\bigskip
\noindent\textit{[Afsnit II--IV: transkription under udarbejdelse.]}

\end{document}
